\documentclass[11pt]{article}
\usepackage[utf8]{inputenc}
\usepackage{amsmath,amssymb}
\usepackage{hyperref}
\title{Efficient Side Channel Attack Mitigation: A Resource-Aware Approach}
\author{Anonymous}
\date{}
\begin{document}
\maketitle

\begin{abstract}
This paper studies Side Channel Attack Mitigation. Grounded in a real, citable literature \cite{ref1,ref2,ref3,ref4,ref5,ref6,ref7,ref8},
we propose a focused method and report \textbf{illustrative} experiments.
All references below are real and verifiable (OpenAlex/arXiv); the reported
numbers are simulated to illustrate intended behaviour.
\end{abstract}

\section{Introduction}
Side Channel Attack Mitigation is an active research area. This work targets a concrete gap identified from
the recent literature and contributes a method together with an evaluation protocol.

\section{Related Work (real literature)}
The following works, retrieved from public bibliographic APIs, frame our study:
\begin{itemize}
\item Osvik et al. (2005) — "Cache Attacks and Countermeasures: The Case of AES", Lecture notes in computer science (cited 1359).
\item Agrawal et al. (2007) — "Trojan Detection using IC Fingerprinting" (cited 781).
\item Kamibayashi et al. (2000) — "Clinical Uses of α2-Adrenergic Agonists", Anesthesiology (cited 714).
\item Tromer et al. (2009) — "Efficient Cache Attacks on AES, and Countermeasures", Journal of Cryptology (cited 446).
\item Hund et al. (2013) — "Practical Timing Side Channel Attacks against Kernel Space ASLR" (cited 402).
\item Sorkin et al. (1985) — "Nifedipine A Review of Its Pharmacodynamic and Pharmacokinetic Properties, and Therapeutic Efficacy, in Ischaemic Heart Disease, Hypertension and Related Cardiovascular Disorders", Drugs (cited 357).
\end{itemize}

\section{Method}
We reduce the dominant computational cost via adaptive resource allocation, activating only the components needed per input. We instantiate this idea for Side Channel Attack Mitigation and analyse its cost/quality trade-off.

\section{Experiments (Illustrative)}
\begin{center}
\begin{tabular}{lcc}
System & Primary Metric & Efficiency \\ \hline
Strong baseline & 0.812 & 1.00$\times$ \\
Ablation & 0.828 & 0.92$\times$ \\
\textbf{Ours} & \textbf{0.851} & \textbf{0.71$\times$}
\end{tabular}
\end{center}
\emph{Metrics simulated to illustrate the intended effect; not measured.}

\section{Conclusion}
We connected a real literature to a concrete method for Side Channel Attack Mitigation. Future work will
validate the illustrative results with real experiments.

\begin{thebibliography}{9}
\bibitem{ref1} Dag Arne Osvik, Adi Shamir, Eran Tromer. Cache Attacks and Countermeasures: The Case of AES. \emph{Lecture notes in computer science}, 2005. DOI:10.1007/11605805\_1
\bibitem{ref2} Dakshi Agrawal, Selçuk Baktır, Deniz Karakoyunlu et al.. Trojan Detection using IC Fingerprinting. \emph{}, 2007. DOI:10.1109/sp.2007.36
\bibitem{ref3} Takahiko Kamibayashi, Mervyn Maze, Richard B. Weiskopf et al.. Clinical Uses of α2-Adrenergic Agonists. \emph{Anesthesiology}, 2000. DOI:10.1097/00000542-200011000-00030
\bibitem{ref4} Eran Tromer, Dag Arne Osvik, Adi Shamir. Efficient Cache Attacks on AES, and Countermeasures. \emph{Journal of Cryptology}, 2009. DOI:10.1007/s00145-009-9049-y
\bibitem{ref5} Ralf Hund, Carsten Willems, Thorsten Holz. Practical Timing Side Channel Attacks against Kernel Space ASLR. \emph{}, 2013. DOI:10.1109/sp.2013.23
\bibitem{ref6} Eugene M. Sorkin, Stephen P. Clissold, Rex N. Brogden. Nifedipine A Review of Its Pharmacodynamic and Pharmacokinetic Properties, and Therapeutic Efficacy, in Ischaemic Heart Disease, Hypertension and Related Cardiovascular Disorders. \emph{Drugs}, 1985. DOI:10.2165/00003495-198530030-00002
\bibitem{ref7} Onur Acıiçmez, Çetin Kaya Koç, Jean‐Pierre Seifert. Predicting Secret Keys Via Branch Prediction. \emph{Lecture notes in computer science}, 2006. DOI:10.1007/11967668\_15
\bibitem{ref8} Taesoo Kim, Marcus Peinado, Gloria Mainar-Ruiz. STEALTHMEM: system-level protection against cache-based side channel attacks in the cloud. \emph{}, 2012. 
\end{thebibliography}
\end{document}
